\chapter{Examples}
\label{sec:examples}

\section{Deterministic Putinar And Full-Box Selection}
\label{sec:full-box-example}

This example selects default Putinar and full-box certificates and an explicit
full-box order without solving an optimization problem.  The reported values
depend only on the public assembly contract, not on incidental YALMIP
decision-variable identifiers.

\begin{lstlisting}[style=dpmatlab]
yalmip('clear')
P = pdvar(1, {[0 1]}, Degree=2);
direct = P >= 0;
putinarDefault = direct.applyPutinar();
boxDefault = direct.applyFullBoxPreorder();
boxHigher = boxDefault.applyFullBoxPreorder(2);
polya = direct.applyPolya(1);
boxFromPolya = polya.applyFullBoxPreorder();
putinarFromBox = boxHigher.applyPutinar();
Fbox = toYalmip(boxDefault);

directState = [direct.UseFullBoxPreorder, direct.FullBoxOrder, ...
    numel(direct.Constraints)]
defaultState = [boxDefault.UseFullBoxPreorder, boxDefault.FullBoxOrder, ...
    numel(boxDefault.Constraints)]
putinarState = [putinarDefault.UsePutinar, putinarDefault.PutinarOrder, ...
    numel(putinarDefault.Constraints)]
higherState = [boxHigher.FullBoxOrder, numel(boxHigher.Constraints)]
replacementState = [boxFromPolya.UsePolya, ...
    boxFromPolya.UseFullBoxPreorder, boxFromPolya.FullBoxOrder]
putinarReplacementState = [putinarFromBox.UsePutinar, ...
    putinarFromBox.UseFullBoxPreorder, putinarFromBox.PutinarOrder]
exportedConstraintCount = length(Fbox)
\end{lstlisting}

\begin{lstlisting}[style=dpoutput]
directState =
     0     0     3
defaultState =
     1     1     5
putinarState =
     1     1     5
higherState =
     2     7
replacementState =
     0     1     1
putinarReplacementState =
     1     0     1
exportedConstraintCount =
     5
\end{lstlisting}

For the degree-two scalar residual, order one uses two PSD Gram blocks followed
by three exact coefficient identities.  Order two matches a degree-four target
and stores two PSD blocks followed by five identities.  The original
\code{direct} object remains unchanged, and the selection made from
\code{polya} replaces P\'olya rather than elevating its already assembled
constraints.  Likewise, selecting Putinar from \code{boxHigher} rebuilds from
the original residual and clears the full-box selection.

\section{Parameter-Dependent Lyapunov Solver Smoke}
\label{sec:parameter-dependent-lyapunov-smoke}

The first solver smoke test in \code{+tests/+pdlmi/test_solver_smoke.m}
contrasts a constant Lyapunov matrix with a degree-one, parameter-dependent
one.  It solves the same decay and positivity constraints in both cases.  The
degree-one solution is then materialized as known \code{pdmat} data and
plotted.  \code{pdvar} is a decision-expression class and does not itself
provide \code{plot}; applying \code{value} to its local YALMIP payloads after
a successful solve supplies the numeric Bernstein coefficients for
\code{pdmat}.

\begin{lstlisting}[style=dpmatlab]
yalmip('clear')
grid = {[0 1]};
A = pdmat(grid, @(x) (1 - x) * [-1 -1; 1 -1] + ...
    x * [-1 -10; 0.1 -1], Degree=1);
solver = 'lmilab';
if exist('mosekopt', 'file') ~= 0
    solver = 'mosek';
end
opts = sdpsettings('solver', solver, 'verbose', 0);

Pc = pdvar(2, grid, Degree=0);
CconstantDecay = Pc * A + A' * Pc <= -1e-10 * eye(2);
CconstantPositive = Pc >= eye(2);
solConstant = optimize([CconstantDecay.toYalmip, ...
    CconstantPositive.toYalmip], [], opts);

Pd = pdvar(2, grid);
CdependentDecay = Pd * A + A' * Pd <= -1e-10 * eye(2);
CdependentPositive = Pd >= eye(2);
solDependent = optimize([CdependentDecay.toYalmip, ...
    CdependentPositive.toYalmip], [], opts);
assert(solDependent.problem == 0, solDependent.info)

Pplot = pdmat(grid, ...
    {cellfun(@value, Pd.LocalValues{1}, UniformOutput=false)}, ...
    Degree=Pd.Degree);
figure(Name="Parameter-dependent Lyapunov matrix")
plot(Pplot, SamplesPerCell=40, LineWidth=1.5)
title("Solved parameter-dependent Lyapunov matrix")
\end{lstlisting}

The figure contains one curve for each entry of the solved \(2\)-by-\(2\)
matrix \(P(\rho)\), sampled on the parameter interval.  It is a diagnostic
visualization of the feasible solution, not a certificate beyond the
coefficient-wise constraints supplied to the solver.  The constant-case result
is retained to exercise the comparison in the smoke test; only a MOSEK solve
is used there to certify its intended infeasibility distinction.

\section{Rate-Dependent Block LMI Solver Smoke}
\label{sec:solver-example}

This example is adapted from \code{+tests/+pdlmi/test_solver_smoke.m}.  The
goal is to find a parameter-dependent \(P(\rho)\) and scalar \(\gamma\) such
that, on every Bernstein coefficient and rate vertex,
\[
\begin{bmatrix}
\dot P + PA + A^\top P & PB & C^\top\\
B^\top P & -\gamma I & D^\top\\
C & D & -\gamma I
\end{bmatrix}
\preceq 0,
\qquad
P(\rho)\succeq 0.
\]
The derivative term \(\dot P\) is built with \code{rhodiff(P, [-1 1])}.  The
comparison operators create \code{pdlmi} wrappers, and \code{toYalmip} hands
the coefficient-wise constraints to YALMIP.  The second listing assigns the
solver- and tolerance-dependent numeric value of \(\gamma\) to
\code{gammaValue}.  For this model, \code{numel(E1.Constraints)} returns 6,
and \code{numel(E2.Constraints)} returns 2.

\begin{lstlisting}[style=dpmatlab]
yalmip('clear')
grid = linspace(0, 1, 2);
A = pdmat(grid, @(x) [-1, 0.5; -1, -2] + ...
    x * [-1.3, -20; 2, -10], Degree=1);
B = pdmat(grid, @(x) [1, -4; -1, -1] + ...
    x * [2.2, 0.5; -6, -5], Degree=1);
Cmat = eye(2);
Dmat = zeros(2);

P = pdvar(2, grid);
diffP = rhodiff(P, [-1 1]);
gamma = pdvar(1, grid, Degree=0);
E1 = [diffP + P * A + A' * P, P * B, Cmat';
      B' * P, -gamma * eye(2), Dmat';
      Cmat, Dmat, -gamma * eye(2)] <= 0;
E2 = P >= 0;
decayConstraintCount = numel(E1.Constraints)
positivityConstraintCount = numel(E2.Constraints)
\end{lstlisting}

\begin{lstlisting}[style=dpoutput]
decayConstraintCount =
     6

positivityConstraintCount =
     2
\end{lstlisting}

\begin{lstlisting}[
  style=dpmatlab,
  title={\textcolor{dpblue}{\sffamily\bfseries\footnotesize Example (continued)}}
]
objective = gamma.LocalValues{1}{1};
solver = 'lmilab';
if exist('mosekopt', 'file') ~= 0
    solver = 'mosek';
end
opts = sdpsettings('solver', solver, 'verbose', 0);
sol = optimize([E1.toYalmip, E2.toYalmip], objective, opts);
assert(sol.problem == 0, sol.info)
gammaValue = value(objective);
assert(isfinite(gammaValue))
gammaValue
\end{lstlisting}

The constraint counts above follow directly from the three independent
assembly axes: physical-cell identity, rate-row identity, and local Bernstein
label.  In this tested example the two wrappers use the same physical grid but
different rate-row and degree structures.

\begin{center}
\begin{tikzpicture}[>=Latex,node distance=7mm and 10mm,
  lane/.style={dp/known,
    font=\scriptsize,align=center,minimum width=64mm,minimum height=12mm,
    inner sep=4pt},
  axes/.style={dp/neutral,
    font=\scriptsize,align=left,minimum width=64mm,minimum height=24mm,
    inner sep=5pt},
  count/.style={dp/contribution,
    font=\scriptsize,align=center,minimum width=64mm,minimum height=12mm,
    inner sep=4pt},
  merge/.style={dp/result,
    font=\scriptsize,align=center,minimum width=75mm,minimum height=12mm,
    inner sep=4pt}]
  \node[lane] (e1) {\code{E1}: rate-dependent block residual\\
    \(R^{(c,v)}[i]\preceq0\)};
  \node[lane,right=of e1] (e2) {\code{E2}: ordinary positivity residual\\
    \(P^{(c)}[i]\succeq0\)};

  \node[axes,below=of e1] (a1) {
    physical cells: \(c\in\{1\}\) \hfill (1)\\
    rate rows: \(v\in\{-,+\}\) \hfill (2)\\
    degree-two local labels: \(i\in\{0,1,2\}\) \hfill (3)};
  \node[axes,below=of e2] (a2) {
    physical cells: \(c\in\{1\}\) \hfill (1)\\
    ordinary row (no rate enumeration) \hfill (1)\\
    degree-one local labels: \(i\in\{0,1\}\) \hfill (2)};

  \node[count,below=of a1] (n1)
    {\(1\times2\times3=6\) finite LMIs};
  \node[count,below=of a2] (n2)
    {\(1\times1\times2=2\) finite LMIs};

  \node[merge,below=14mm of n1,xshift=37mm] (export)
    {\code{E1.toYalmip} and \code{E2.toYalmip}\\
     enter one finite YALMIP model};
  \node[merge,below=of export] (solve)
    {\code{optimize([E1.toYalmip, E2.toYalmip], objective, opts)}};

  \draw[->,thick,dpblue] (e1)--(a1);
  \draw[->,thick,dpblue] (e2)--(a2);
  \draw[->,thick,dpblue] (a1)--(n1);
  \draw[->,thick,dpblue] (a2)--(n2);
  \draw[->,thick,dpblue] (n1)--(export);
  \draw[->,thick,dpblue] (n2)--(export);
  \draw[->,thick,dpblue] (export)--(solve);
\end{tikzpicture}
\end{center}

The numbers (6) and (2) are certificate-assembly counts: they state how
many finite semidefinite constraints represent \code{E1} and \code{E2}.
They are not evidence that the optimization succeeds.  Feasibility is checked
separately by \code{sol.problem}, and the finite objective is checked after the
shared \code{optimize} call.
